\SourceLeaf{401}%

\pdfbookmark[1]{II. Bestimmung der nächsten ganzen Zahl}{complexe-II}
\begin{center}[II.]\par\medskip
BESTIMMUNG DER NACHSTEN GANZEN ZAHL.
\end{center}
Es sei
\[
\begin{aligned}
\varepsilon^3&=1,\quad m=a+b\varepsilon+c\varepsilon\varepsilon\\
2a-b-c&=A+\alpha\\
2b-c-a&=B+\beta\\
2c-a-b&=C+\gamma
\end{aligned}
\]
wo $A,B,C$ ganze Zahlen; $\alpha,\beta,\gamma$ positive echte Brüche sind. Man hat dann
\[
A+B+C+\alpha+\beta+\gamma=0
\]
also drei Fälle zu unterscheiden:

I. $\alpha+\beta+\gamma=0$, folglich $\alpha=0$, $\beta=0$, $\gamma=0$

1, $A\equiv B\equiv C\ (\text{mod. }3)$. Hier ist $m$ selbst eine ganze Zahl.

2, $A-B\equiv B-C\equiv C-A\equiv\pm1\ (\text{mod. }3)$. Hier ist $m\pm\frac{\varepsilon-\varepsilon\varepsilon}{3}\pdot\varepsilon^n$ eine ganze Zahl.

II. $\alpha+\beta+\gamma=1$.

Hier ist
\[
\begin{aligned}
&A+B\varepsilon+C\varepsilon\varepsilon+1\\
&A+B\varepsilon+C\varepsilon\varepsilon+\varepsilon\\
&A+B\varepsilon+C\varepsilon\varepsilon+\varepsilon\varepsilon
\end{aligned}
\]
\SourceLeaf{402}%
jedes durch $1-\varepsilon$ theilbar, und eine dieser Zahlen durch $3$. Der Quotient oder
\[
m+\frac{\varepsilon^n-\alpha-\beta\varepsilon-\gamma\varepsilon\varepsilon}{3}
\]
die gesuchte ganze Zahl.

III. $\alpha+\beta+\gamma=2$

Hier sind
\[
\begin{aligned}
&A+B\varepsilon+C\varepsilon\varepsilon+\varepsilon+\varepsilon\varepsilon\\
&A+B\varepsilon+C\varepsilon\varepsilon+\varepsilon\varepsilon+1\\
&A+B\varepsilon+C\varepsilon\varepsilon+1+\varepsilon
\end{aligned}
\]
durch $1-\varepsilon$ und eine dieser Zahlen durch $3$ theilbar. Der Quotient, oder
\[
m+\frac{\varepsilon^n(\varepsilon+\varepsilon\varepsilon)-\alpha-\beta\varepsilon-\gamma\varepsilon\varepsilon}{3}
\]
ist die gesuchte ganze Zahl.

In allen drei Fällen hat der Rest die Form
\[
x+y\varepsilon+z\varepsilon\varepsilon
\]
so dass $x,y,z$ ohne Rücksicht auf das Zeichen kleiner als $\frac13$ und $x+y+z=0$ wird. Dadurch wird aber nothwendig
\[
xx+yy+zz=2xx-2yz=2yy-2xz=2zz-2xy<\tfrac29
\]
weil von den drei Grössen $x,y,z$ nothwendig zwei einerlei Zeichen haben. Folglich ist der Determinant des Restes
\[
=\tfrac32(xx+yy+zz)<\tfrac13\qquad\text{Q. E. D.}
\]

Die Bestimmung der nächsten ganzen Zahl geschieht \emph{bequemer} auf folgende Art. Es sei vorgegeben $a+b\varepsilon+c\varepsilon\varepsilon=m$, man setze
\[
\begin{aligned}
b-a&=C+\gamma\\
c-b&=A+\alpha\\
a-c&=B+\beta
\end{aligned}
\]
wo $A,B,C$ die nächst kleinern ganzen Zahlen; $\alpha,\beta,\gamma$ positive Brüche sind.
Hier sind drei Fälle zu unterscheiden:

\SourceLeaf{403}%
I. $\alpha+\beta+\gamma=0$, so ist $m$ selbst ganze Zahl

II. $\alpha+\beta+\gamma=1$, so ist die nächste ganze Zahl
\[
\begin{array}{ll}
B+(B+C)\varepsilon&\text{wenn }\alpha\text{ der grösste Bruch ist.}\\
\qquad C\varepsilon+(A+C)\varepsilon\varepsilon&\qquad\beta\\
A+B\quad\pdot\quad+A\varepsilon\varepsilon&\qquad\gamma
\end{array}
\]

II. $\alpha+\beta+\gamma=2$, so ist die nächste ganze Zahl
\[
\begin{array}{ll}
B+1+(B+C+2)\varepsilon&\text{wenn }\alpha\text{ der kleinste Bruch ist.}\\
\qquad(C+1)\varepsilon+(A+C+2)\varepsilon\varepsilon&\qquad\beta\\
A+B+2\quad\pdot\quad+(A+1)\varepsilon\varepsilon&\qquad\gamma
\end{array}
\]
In II, $1$ ist der Rest $\beta+(\beta+\gamma)\varepsilon$, dessen Determinant
\[
\begin{aligned}
&=\beta\beta+\beta\gamma+\gamma\gamma\\
&=\tfrac13-\tfrac16[(\alpha-\beta)(1+3\beta)+(\alpha-\gamma)(1+3\gamma)]
\end{aligned}
\]

\emph{Noch einfacher so:}

Man ordne die Brüche $a-[a],b-[b],c-[c]$ nach ihrer Grösse: so heissen sie der Reihe nach $p,q,r$. Sind alle drei gleich gross, so ist $m$ eine ganze Zahl. Sind sie aber ungleich, so sei $t$ ein beliebiger Bruch zwischen
\[
\begin{array}{llll}
p\text{ und }q,&\text{jenachdem}&q-p&\text{am grössten ist}\\
q\text{ und }r&&r-q&\\
r\text{ und }1+p&&1+p-r&
\end{array}
\]
Sodann ist
\[
[a-t]+[b-t]\varepsilon+[c-t]\varepsilon\varepsilon
\]
die nächste ganze Zahl.
\BookRule{.3\textwidth}

\SourceLeaf{404}%
\pdfbookmark[1]{III.}{complexe-III}
\begin{center}[III.]\end{center}
Es sei $\varepsilon^5=1$
\[
\begin{aligned}
a+b\varepsilon+c\varepsilon\varepsilon+d\varepsilon^2+e\varepsilon^4&=q'\\
a+b\varepsilon^{-1}+c\varepsilon^{-2}+d\varepsilon^{-3}+e\varepsilon^{-4}&=q''''
\end{aligned}
\]
\[
\begin{aligned}
(a-b)^2+(b-c)^2+(c-d)^2+(d-e)^2+(e-a)^2&=2p'\\
(a-c)^2+(b-d)^2+(c-e)^2+(d-a)^2+(e-b)^2&=2p''
\end{aligned}
\]
\[
\begin{aligned}
q'q''''&=-p'\varepsilon-p''\varepsilon\varepsilon-p''\varepsilon^3-p'\varepsilon^4=P'\\
q''q'''&=-p'\varepsilon\varepsilon-p''\varepsilon^4-p''\varepsilon-p'\varepsilon^3=P''
\end{aligned}
\]

\[
\text{Determinant}=P'P''=-p'p'+3p'p''-p''p''
\]
\[
\begin{aligned}
\text{Mensura}&=2p'+2p''=2P'+2P''\\
&=5(aa+bb+cc+dd+ee)-(a+b+c+d+e)^2
\end{aligned}
\]
Multiplicando per $1-\varepsilon$ fit mensura nova $=8p'$
\[
\text{Höchste Mensur}=2\left(\frac{\sin72^\circ}{\sin36^\circ}+\frac{\sin36^\circ}{\sin72^\circ}\right)\surd D=4{,}472\surd D
\]

\emph{Modulus} $=1-\varepsilon$
\[
\begin{aligned}
1-\varepsilon&=x\\
\varepsilon&=1-x\\
\varepsilon\varepsilon&=1-2x+xx\\
\varepsilon^3&=1-3x+3xx-x^3\\
\varepsilon^4&=1-4x+6xx-4x^3+x^4\\
&=-4+6x-4xx+x^3
\end{aligned}
\]
Also
\[
\begin{aligned}
\frac{1-\varepsilon^n}{1-\varepsilon}&\equiv n\text{ mod. }(1-\varepsilon)\\
\varepsilon^n&\equiv1-nx\text{ mod. }(1-\varepsilon)^2\\
\left(\frac{\varepsilon+\varepsilon^4}{\varepsilon\varepsilon+\varepsilon^3}\right)^n&\equiv1+nxx\text{ mod. }(1-\varepsilon)^3
\end{aligned}
\]
Also eine Zahl, welche $\equiv1\text{ mod. }(1-\varepsilon)^3$ kann \emph{nur} dann eine Einzahl sein, wenn sie zugleich $\equiv1\ (\text{mod. }5)$.
\BookRule{.3\textwidth}

\SourceLeaf{405}%
\pdfbookmark[1]{IV. Einiges über die Mensur der Zahlen}{complexe-IV}
\begin{center}[IV.]\par\medskip
EINIGES UBER DIE MENSUR DER ZAHLEN.
\end{center}
Es sei $\varepsilon^n=1$, $n$ Primzahl
\[
\begin{aligned}
m&=a+a'\varepsilon+a''\varepsilon\varepsilon+a'''\varepsilon^3+\ldots+a^{(n-1)}\varepsilon^{n-1}=f\varepsilon\\
D&=f\varepsilon\pdot f\varepsilon\varepsilon\pdot f\varepsilon^3\ldots f\varepsilon^{n-1}\\
f\varepsilon\pdot f\varepsilon^{n-1}&=-b'(\varepsilon+\varepsilon^{n-1})-b''(\varepsilon^2+\varepsilon^{-2})\\
&\qquad-b'''(\varepsilon^3+\varepsilon^{-3})\ldots
\end{aligned}
\]
so ist
\[
\begin{aligned}
2b'&=(a-a')^2+(a'-a'')^2+(a''-a''')^2+\text{etc.}\\
2b''&=(a-a'')^2+(a'-a''')^2+(a''-a'''')^2+\text{etc.}\\
&\text{etc.}
\end{aligned}
\]

hier sind also $b',b'',b'''\dotstwo$ lauter positive Grössen; sie heissen \emph{Partialmensuren} von $m$, so wie ihre Summe
\[
\begin{aligned}
b'+b''+b'''+\text{etc.}={}&n(aa+a'a'+a''a''+\dotstwo)\\
&\qquad-(a+a'+a''+\text{etc.})^2
\end{aligned}
\]
die \emph{Generalmensur}. Setzt man
\[
f\varepsilon\pdot f\varepsilon^{n-1}=c',\quad f\varepsilon\varepsilon\pdot f\varepsilon^{n-2}=c''\quad\text{etc.}
\]
so ist
\[
\begin{aligned}
&c'+c''+c'''+\text{etc.}+c^{\frac12(n-1)}\\
&\qquad=b'+b''+b'''+\text{etc.}+b^{\frac12(n-1)}
\end{aligned}
\]
\[
\begin{aligned}
&c'(\varepsilon+\varepsilon^{n-1})+c''(\varepsilon\varepsilon+\varepsilon^{n-2})+c'''(\varepsilon^3+\varepsilon^{n-3})+\text{etc.}\\
&\qquad=2(b'+b''+b'''+\text{etc.}+b^{\frac12(n-1)})-nb'
\end{aligned}
\]
\[
\begin{aligned}
&c'(2-\varepsilon-\varepsilon^{n-1})+c''(2-\varepsilon\varepsilon-\varepsilon^{n-2})\\
&\qquad+c'''(2-\varepsilon^3-\varepsilon^{n-3})+\text{etc.}=nb'
\end{aligned}
\]
\[
b'>\frac{n-1}{2n}(nD)^{\frac2{n-1}},\quad b'+b''+b'''+\text{etc.}>\frac{n-1}{2}\pdot D^{\frac2{n-1}}
\]

Ist allgemein
\[
f\varepsilon\pdot f\varepsilon^{n-1}=A+A'\varepsilon+A''\varepsilon\varepsilon+A'''\varepsilon^3+\dotstwo
\]
so ist die Generalmensur $\triangle=-A-A'-A''-\text{etc.}+nA$

Mensur von $(1+\varepsilon)f\varepsilon\ldots$
\[
\triangle'=4\triangle-2n(A-A')=4\triangle-2nb'
\]
Ist $a+a'+a''+\dotstwo=0$, so ist $\triangle=n(aa+a'a'+a''a''+\text{etc.})$

und ist $A+A'+A''+\ldots=0$, so ist $\triangle=nA$, $\triangle'=n(2A+2A')$

Ist also einer der Coëfficienten $A',A''$ etc. negativ und absolut grösser als $\frac12A$, so lässt sich die Mensur salvo determinante herabbringen.
\BookRule{.3\textwidth}

\SourceLeaf{406}%
\pdfbookmark[1]{V.}{complexe-V}
\begin{center}[V.]\end{center}
Sollte sich bestätigen, dass jede Einheitszahl blos aus Factoren von der Form
\[
\frac{\varepsilon^\alpha-\varepsilon^\beta}{\varepsilon^\gamma-\varepsilon^\delta}
\]
zusammengesetzt wäre, so würde folgender Satz bewiesen sein:

\emph{Ist $f(\varepsilon)$ eine Einheitszahl, so ist}
\[
\frac{f(\varepsilon)}{f(\varepsilon^{-1})}=\varepsilon^n
\]

Auch ohne \emph{jenen} Satz \emph{vorauszusetzen} ist der Schlusssatz leicht zu beweisen.
Es sei
\[
\frac{f\varepsilon}{f\varepsilon^{-1}}=F\varepsilon
\]
so ist
\[
F\varepsilon\pdot F\varepsilon^{-1}=1
\]
woraus mit Hülfe der Lehre von der Mensur leicht gefolgert wird, dass
\[
F\varepsilon=\pm\varepsilon^n
\]
Das untere Zeichen ist aber unmöglich, weil sonst $f\varepsilon$ durch $1-\varepsilon$ theilbar sein müsste.

Dass der Determinant einer von $0$ verschiedenen Zahl nicht $=0$ sein könne, lässt sich leicht beweisen. Wenn der Determinant durch $m$ theilbar ist, so ist die Zahl selbst durch $1-\varepsilon$ theilbar; folglich wenn der Determinant durch $m^{m-1}$ theilbar ist, muss die Zahl selbst durch $m$ theilbar sein. Welches absurd ist, da beim Det. $0$ die Zahl erst salvo Det. so oft durch $m$ dividirt werden könnte, bis sie nicht mehr theilbar wäre. Der erste Satz aber erhellt so. Es sei die vorgegebene Zahl
\[
a+b\varepsilon+c\varepsilon\varepsilon+\text{etc.}\equiv a+b+c\dotstwo\text{ mod. }1-\varepsilon
\]
also Determinans $\equiv(a+b+c\dotstwo)^{m-1}\text{ mod. }1-\varepsilon$.
\BookRule{.3\textwidth}

\SourceLeaf{407}%
\pdfbookmark[1]{VI.}{complexe-VI}
\begin{center}[VI.]\end{center}
Es sei $\varepsilon^n=1$
\[
\begin{aligned}
f\varepsilon&=a+b\varepsilon+c\varepsilon\varepsilon+d\varepsilon^3+\text{etc.}\\
m&=\text{Determinans dieser Zahl}\\
\frac{m}{f\varepsilon}&=f\varepsilon\varepsilon\pdot f\varepsilon^3\ldots f\varepsilon^{n-1}\\
&=A+B\varepsilon+C\varepsilon\varepsilon+\text{etc.}=F\varepsilon
\end{aligned}
\]
Der Zahl $f\varepsilon$ entspricht eine Wurzel der Congruenz $x^n\equiv1\ (\text{mod. }m.)$ Es sei dieselbe $r$. Man hat
\[
\begin{aligned}
nA&=F1+F\varepsilon+F\varepsilon\varepsilon+\dotstwo\\
nB&=F1+\varepsilon^{-1}F\varepsilon+\varepsilon^{-2}F\varepsilon\varepsilon+\dotstwo\\
nC&=F1+\varepsilon^{-2}F\varepsilon+\varepsilon^{-4}F\varepsilon\varepsilon+\dotstwo\\
&\text{etc.}
\end{aligned}
\]

also, da $F\varepsilon\varepsilon,F\varepsilon^3,F\varepsilon^4$ etc. durch $f\varepsilon$ theilbar sind,
\[
\begin{gathered}
nA-F1-\varepsilon(nB-F1)\\
nA-F1-\varepsilon\varepsilon(nC-F1)\\
nA-F1-\varepsilon^3(nD-F1)\\
\text{etc.}
\end{gathered}
\]
alle durch $f\varepsilon$ theilbar, oder auch
\[
\begin{gathered}
n(A-B)-\varepsilon n(B-C)\\
n(B-C)-\varepsilon n(C-D)\\
\text{etc.}
\end{gathered}
\]
durch $f\varepsilon$ theilbar; folglich [wenn $f\varepsilon$ durch $1-\varepsilon$, und $F\varepsilon$ durch eine ganze reelle Zahl nicht theilbar ist]
\[
\varepsilon\equiv\frac{A-B}{B-C}\equiv\frac{B-C}{C-D}\equiv\frac{C-D}{D-E}\text{ etc. }(\text{mod. }f\varepsilon)
\]
\BookRule{.3\textwidth}

\SourceLeaf{408}%
\pdfbookmark[1]{Bemerkungen zu den complexen Zahlen}{complexe-Bemerkungen}
\begin{center}BEMERKUNGEN.\end{center}
\BookRule{.15\textwidth}
Die hier unter der gemeinsamen Ueberschrift, zur Theorie der complexen Zahlen, zusammengestellten Untersuchungen bilden zerstreute Notizen in der Handschrift. Sie enthalten die wesentlichen Momente des Beweises vom \textsc{Fermat}schen Satze für die dritte und fünfte Potenz. Die aus dritten Wurzeln der Einheit zusammengesetzten Zahlen sind in unvollständigen hier nicht abgedruckten Aufzeichnungen sowohl mit Hülfe der Theorie der binären quadratischen Formen, als auch der Kreistheilung untersucht. Bei Gelegenheit der Anwendung der letztern und zwar während der Ausarbeitung der Abhandlung \emph{Disquisitionum circa aequationes puras ulterior evolutio} ist noch die ternäre cubische Form aufgestellt, in welche $27\frac{x^n-1}{x-1}$ für eine Primzahl $n\equiv1\text{ mod. }3$ verwandelt werden kann, und zugleich die Theorie der Composition der mit jener verwandten Form $X^3+mY^3+mmZ^3-3mXYZ$ entwickelt.

Die in den Untersuchungen des Bruchstück [I] vorausgesetzte Eigenschaft der aus dritten Wurzeln der Einheit gebildeten ganzen Zahlen, dass jede nur auf Eine Weise in Primfactoren zerlegt werden kann, ergibt sich aus dem \textsc{Euclid}ischen Verfahren, die gemeinsamen Theiler zweier Zahlen zu bestimmen, wenn dabei der unter [II] abgeleitete Satz über die nächste ganze Zahl für irgend eine vorgegebene Bruchzahl in Anwendung gebracht wird.

Dass dieselbe Fundamentaleigenschaft auch den aus fünften Wurzeln der Einheit zusammengesetzten Zahlen zukommt, folgt daraus, dass der nach einer ganz analogen Regel wie in [II] gebildete Bruchrest entweder von $m$ oder doch von $m$ multiplicirt in eine geeignete Einheitszahl $E$ so beschaffen ist, dass er durch Subtraction von der vorgegebenen Zahl $mE$ eine ganze Zahl entstehen lässt und dass sein Determinant die Einheit nicht übertrifft. Die Einheitszahlen lassen sich aber, wie in [III] angedeutet, aus der Theorie der binären quadratischen Formen vom Determinant $5$ in Verbindung mit der Zerlegung irgend einer reellen Primzahl in vier Factoren (z. B. $11=\text{Det. }(2+\varepsilon)$) ableiten, nemlich als Producte der Potenzen von $\varepsilon$ und $1+\varepsilon$.
\begin{flushright}\textsc{Schering.}\end{flushright}
\BookRule{.3\textwidth}

